\documentclass[10pt,a4paper]{amsart}
\usepackage[margin=19mm]{geometry}
\usepackage{amsmath,amssymb,mathtools}
\usepackage[scr=rsfs,cal=euler]{mathalfa}
\usepackage{xfrac}
\usepackage[unicode,hidelinks,bookmarks=false]{hyperref}
\newcommand{\ie}{i.e.\ }
\newcommand{\cf}{cf.\ }
\newcommand{\resp}{respectively\ }
\newcommand{\Subsection}{\subsection}
\providecommand{\nobreakdash}{\nobreak\hbox{-}}
\setlength{\emergencystretch}{2em}
\multlinegap0pt
\usepackage{amssymb}
\usepackage{mathtools}
\usepackage{bm}
\usepackage{float}
\usepackage[shortlabels]{enumitem}
\usepackage{siunitx}
\usepackage{xcolor}
\newtheorem{theorem}{Theorem}
\newcommand{\calB}{\mathcal B}
\newcommand{\calI}{\mathcal I}
\newcommand{\dd}{\mathrm{d}}
\title{Quantifying Side-Channel Leakage in Public Metrology Releases\\[1.5ex]{\large\itshape Screened EUV Roughness Spectra as a Case Study}}
\author{Faruk Alpay, Taylan Alpay}
\date{Source: arXiv:2606.02934v1 (CC BY 4.0)}
\begin{document}
\maketitle

\noindent\emph{Source and licence.} Quantifying Side-Channel Leakage in Public Metrology Releases\\[1.5ex]{\large\itshape Screened EUV Roughness Spectra as a Case Study}. Version arXiv:2606.02934v1 (CC BY 4.0), distributed by arXiv under the Creative Commons Attribution 4.0 International licence, http://creativecommons.org/licenses/by/4.0/, as recorded in the arXiv licence metadata for this exact version and shown on the abstract page https://arxiv.org/abs/2606.02934v1. Full licence notice: \texttt{licenses/arxiv-2606.02934-CC-BY-4.0.txt}.\par\smallskip

\noindent\textit{Editorial scope.} Counted unit: the article's theorem thm:k9 (source0.tex lines 408--416). The article's complete original proof of it is delivered with it (source0.tex lines 418--459, one \texttt{proof} environment of 42 lines). No mathematics is written, repaired or re-derived: the statement and the proof are the article's own lines, in the article's own notation, and no retained line is substituted or re-flowed.  No further source-local context is needed: every cross-reference of every retained line resolves inside the two delivered spans.  Nothing was omitted from any retained span, so the package carries no omission record.  No retained line contains a citation, so the wrapper carries no bibliography and no bibliography entry.  No retained line contains a figure environment, an image-inclusion command or drawing code, so no image is delivered and the package declares no asset dependency.  Editorial formatting only, in addition: an AMS \texttt{amsart} preamble that restates the article's own declarations and macros used by the retained lines, namely the environments \texttt{theorem} on separate counters, together with the standard packages those lines need.  The retained environments print in this document as Theorem 1 in order of appearance, whereas the article prints the same items with the numbering of its own full document.  The complete native source of the article as distributed in the arXiv e-print for this exact version is bundled unchanged as \texttt{sources/201943/source0.tex}.
\begin{theorem}[Transport-knee leakage exponent]
\label{thm:k9}
For flat \(w\) and \(K\lambda\ll1\),
\begin{equation}
\calI_{\lambda\mid\alpha,\beta}(K)=\frac{64}{1225}\,w\lambda^6K^9+O(w\lambda^8K^{11}),
\label{eq:k9}
\end{equation}
so for \(\lambda_\pm=\lambda\pm\Delta\lambda/2\), \(C_R=\tfrac18(\Delta\lambda)^2\calI_{\lambda\mid\alpha,\beta}(K)+O((\Delta\lambda)^3)\).
\end{theorem}

\begin{proof}
Let \(\calB=\operatorname{span}\{1,k^2\}\) and \(R_\calB=I-\Pi_{\alpha,\beta}\). For \(K\lambda\ll1\),
\[
g_\lambda(k)=-\frac{3\lambda k^2}{1+\lambda^2k^2}
=-3\lambda k^2+3\lambda^3k^4+O(\lambda^5k^6).
\]
Since \(k^2\in\calB\), the first term is annihilated by \(R_\calB\), and therefore
\[
R_\calB g_\lambda=3\lambda^3R_\calB k^4+O(\lambda^5R_\calB k^6).
\]
The cross term between these two displayed terms is \(O(w\lambda^8K^{11})\), because \(\|R_\calB k^4\|^2=O(K^9)\) and \(\|R_\calB k^6\|^2=O(K^{13})\); the square of the remainder is of higher order. It remains to compute \(\|R_\calB k^4\|^2\).

Write the projection of \(k^4\) onto \(\calB\) as \(c_0+c_2k^2\). With \(\mu_j=\int_0^Kk^j\,\dd k=K^{j+1}/(j+1)\), the normal equations are
\[
\begin{pmatrix}
\mu_0&\mu_2\\
\mu_2&\mu_4
\end{pmatrix}
\begin{pmatrix}c_0\\c_2\end{pmatrix}
=
\begin{pmatrix}\mu_4\\\mu_6\end{pmatrix}.
\]
The determinant is \(K^6(1/5-1/9)=4K^6/45\), so Cramer's rule gives
\[
c_0=-\frac{3}{35}K^4,\qquad c_2=\frac67K^2 .
\]
By orthogonality,
\[
\|R_\calB k^4\|^2
=\int_0^Kk^8\,\dd k-c_0\int_0^Kk^4\,\dd k-c_2\int_0^Kk^6\,\dd k
=\left(\frac19+\frac{3}{175}-\frac{6}{49}\right)K^9
=\frac{64}{11025}K^9 .
\]
Consequently
\[
\calI_{\lambda\mid\alpha,\beta}(K)
=w\|R_\calB g_\lambda\|^2
=9\lambda^6w\,\frac{64}{11025}K^9+O(w\lambda^8K^{11})
=\frac{64}{1225}w\lambda^6K^9+O(w\lambda^8K^{11}).
\]
For the pair \(\lambda_\pm=\lambda\pm\Delta\lambda/2\), the local Chernoff expansion of a regular one-parameter likelihood after Schur elimination is \(C_R=\tfrac18(\Delta\lambda)^2\calI_{\lambda\mid\alpha,\beta}(K)+O((\Delta\lambda)^3)\), which gives the final statement.
\end{proof}

\end{document}
